\documentclass[11pt]{article}
\usepackage[margin=1.1in]{geometry}
\usepackage{amsmath, amssymb, amsthm}
\usepackage{booktabs}
\usepackage{xcolor}
\definecolor{linknavy}{RGB}{25, 55, 125}
\usepackage[colorlinks=true, linkcolor=linknavy, citecolor=linknavy, urlcolor=linknavy]{hyperref}

\newtheorem{theorem}{Theorem}[section]
\newtheorem{proposition}[theorem]{Proposition}
\newtheorem{lemma}[theorem]{Lemma}
\newtheorem{conjecture}[theorem]{Conjecture}
\theoremstyle{remark}
\newtheorem{remark}[theorem]{Remark}

\newcommand{\C}{\mathbb{C}}
\newcommand{\Q}{\mathbb{Q}}
\newcommand{\Gm}{\mathbb{G}_m}

\title{Planar Keller Maps after the Counterexample:\\
Machine-Certified Newton-Polygon Exclusions,\\
the Staircase Transport, and the Recalibrated Frontier}
\author{Felipe Santiba\~nez-Leal\\
\small CAOS Research program, \texttt{github.com/fsantibanezleal/CAOS\_RESEARCH}}
\date{2026-07-22 (working manuscript, v0.02; split from the foundational companion)}

\begin{document}
\maketitle

\begin{abstract}
The July 2026 counterexample settled the Jacobian conjecture in dimensions $N \ge 3$ and left
$JC(2)$ as the open case. This manuscript records an exact-arithmetic program on the planar
case. The Keller condition $J(P, Q) = 1$ is linear in $Q$ for fixed $P$; the resulting machine
(eliminate, parametrize the consistency variety, complete, invert) closes the columns
$\min\deg \le 2$ uniformly (one closed inverse for the whole quasi-triangular class), forces
perfect-power tops at bidegrees $(2, n)$ and $(3, n)$, and inverts every consistent stratum it
has met. Four unconditional exclusion theorems follow from a weight-class calculus: a two-term
$P = x + a x^u y^v$ is never a Keller component at any partner degree; neither is any
lower-weight perturbation of it; no Keller component carries a Newton top edge anchored at the
linear vertex $x$ with interior direction; and if $x$ is a vertex of the Newton polygon then
$P = x + f(y)$ exactly (the vertex dichotomy). The residual open core is the transport of the
Keller constant along the staircase of Newton vertices swallowing $x$: we show the completion
window is block-triangular under the edge grading, exhibit the classwise transport chain as an
executable elimination whose obstruction always localizes at the constant's class, and prove the corresponding
fifth exclusion: for primitive tops the certificate-tower induction closes it
UNCONDITIONALLY at all partner degrees (cleared base certificates with monomial pairings
such as $-13860\,a^7$, plus a resonance-reduction lemma showing window growth can never
repair the constant's row). A source-audited recalibration places the program
honestly: our composite-gcd certificates replicate literature-covered territory (the verified
floor is $\gcd \le 8$, primes, $2p$, then Heitmann's $B \ge 16$ and the
Guccione--Guccione--Valqui dichotomy $B = 16$ or $B > 20$), and the live frontier is the
$B = 16$ normal form together with the single surviving degree pair $(72, 108)$ below $125$.
A new multiplication formula for $x^m$-anchored edges shows the $B = 16$ leading form cannot
produce the Keller constant on its top edge for $m \ge 2$: that case, too, is a staircase
transport. All computations are exact and reproducible; every claim is labeled machine-verified
[MV], derived [D], or conjectural [C].
\end{abstract}

\section{Introduction}

Let $P, Q \in \C[x, y]$ with $J(P, Q) := P_x Q_y - P_y Q_x \in \C^\times$ (a \emph{Keller
pair}; each member is a \emph{Keller component}). The two-variable Jacobian conjecture $JC(2)$
asserts every Keller pair is an automorphism pair. Following the July 2026 counterexample in
dimension 3 \cite{Alpoge2026}, $JC(2)$ is the surviving case, and the classical planar theory
(Newton polygons, approximate roots, gcd coverage) is its context. This manuscript is the
planar half of a two-manuscript record: the foundational companion documents the
three-dimensional aftermath (validation, the seed family, the geometry of non-properness,
two-dimensional equivariant rigidity); this one documents the planar program. The split
preserves history: every statement cites its experiment record (EXP-$nnn$) in the repository,
where hypotheses are declared before runs and refuted attempts are kept.

Labels: \textbf{[MV]} machine-verified exact identity; \textbf{[D]} derived analytically and
certified on instances; \textbf{[C]} conjecture. Machine-verification is exact sympy
arithmetic over $\Q$ or $\Q(i)$; no floats anywhere.

Two structural inputs from the companion are used freely. First, every $\Gm$-equivariant
Keller map of $\C^2$ is linear (equivariant rigidity, EXP-010), so nothing planar is
equivariant beyond the affine group. Second, at every weight ray the leading pair of a planar
Keller map is Jacobian-dependent (the tropical bridge, EXP-013), which is the entry point of
all edge analyses below.

\section{The planar machine}

The Keller condition is bilinear in the coefficient vectors of $P$ and $Q$, hence
\emph{linear in $Q$ given $P$}. In the affine gauge (identity linear parts) this yields a
machine per degree pair: eliminate $Q$ (the consistency ideal in $P$'s coefficients),
parametrize, complete linearly, and invert or fiber-test (EXP-017, EXP-019, EXP-020).

\begin{theorem}[{[MV]}; uniform closure at $\min\deg \le 2$ (EXP-021/022)]\label{thm:uniform}
Write $\ell = \alpha x + \beta y$, $\beta \ne 0$. In shear coordinates the Keller condition
for $P = x + f(\ell)$ (any polynomial $f$) is a linear first-order PDE with characteristic
invariant $P$; its complete polynomial solution space is $Q = \ell/\beta + H(P)$, $H$
arbitrary, and the pair inverts by one closed formula. Consequently every planar Keller map
with $\min(\deg P, \deg Q) \le 2$ is invertible, and the whole quasi-triangular class
$P \sim x + f(\ell)$ is invertible at every degree.
\end{theorem}

\begin{theorem}[{[MV]}; forced tops at $(2, n)$ and $(3, n)$ (EXP-019/020/022/039)]
\label{thm:forced}
The $(2,3)$ consistency ideal is the single equation $(4A_0A_2 - A_1^2)^2 = 0$ ($P$'s top
form is a perfect square); the $(2,4)$ elimination returns the same ideal. At $(3, 4)$, the
$6 \times 5$ linear map $Q_4 \mapsto J(P_3, Q_4)$ is rank-deficient exactly when $P_3$ is a
perfect cube: all $5 \times 5$ minors vanish on the cube parametrization, every minor lies in
the Hessian-covariant ideal, and completeness follows from $GL_2$-equivariance
($J(f \circ A, g \circ A) = \det(A)\, J(f, g) \circ A$, verified identically) plus full rank
at the two non-cube orbit representatives $x^2 y$ and $x y (x + y)$. On the cube stratum,
sampled completions are Keller and invert with explicit polynomial inverses ($4/4$ by lex
Gr\"obner extraction). At $(3,3)$ the consistency ideal forces the quasi-triangular alignment
at the radical level.
\end{theorem}

The degree columns close by classical coverage: $\gcd(2, n) \le 2$ and $\gcd(3, n) \in
\{1, 3\}$ are inside the verified $\gcd$ landscape of Section~\ref{sec:frontier}, so the
machine's contribution there is the \emph{constructive} loop (explicit consistency varieties
and inverses), stated as replication. The classical degree descent is an algorithm on this
machine: subtract $c P^k$ while the tops are power-proportional, finish with
Theorem~\ref{thm:uniform}; it explicitly inverted a deterministic library of composed Keller
maps, and rotate-descent (rotation of linear-base tops plus de Jonquieres subtraction,
EXP-034) linearizes that entire library in $0$ to $6$ steps. Its exact limit: mixed-corner
tops stay mixed under rotation; those are the classical hard shapes, and they are where the
exclusion theorems below take over.

\section{Four unconditional exclusions}

Throughout, $P$ has linear part $x$ after affine gauge. The weight-class calculus: under
$w = (v, 1 - u)$, the operator $J(P, \cdot)$ shifts weight by a constant, so the Keller
equation decouples into weight classes; the technique goes back to Dixmier and Joseph
\cite{Joseph1975} and is standard in the polygon literature \cite{GGV2017}; the statements
below appear to be new in the stated generality (Section~\ref{sec:positioning}).

\begin{theorem}[{[MV/D]}; the weight-class theorem (EXP-029)]\label{thm:t1}
For $u \ge 2$, $v \ge 1$, $a \ne 0$, the polynomial $P = x + a\,x^u y^v$ is not a Keller
component: no $Q$, of any degree, satisfies $J(P, Q) \in \C^\times$.
\end{theorem}

\begin{proof}[Proof shape]
$P$ is $(v, 1-u)$-homogeneous; the constant is reachable only from the class of $y$, spanned
by the ray $g_s = x^{ks} y^{ms+1}$ ($k = (u-1)/d$, $m = v/d$, $d = \gcd(u-1, v)$); there the
operator is banded, $J(P, g_s) = (ms+1)\,e_s + a B_s\,e_{s+d}$ with $B_s$ a positive integer,
and the chain recursion contradicts the last row of every truncation.
\end{proof}

\begin{theorem}[{[MV/D]}; lower-weight perturbations never rescue (EXP-030/031)]\label{thm:t2}
With $u, v, a$ as above and $R$ any polynomial whose monomials have degree $\ge 2$ and
$(v, 1-u)$-weight strictly below $v$: $P = x + a\,x^u y^v + R$ is not a Keller component.
\end{theorem}

\begin{theorem}[{[MV/D]}; every $x$-anchored edge falls (EXP-032)]\label{thm:t3}
Let $z = x^k y^m$, $k, m \ge 1$, and $E = x\,\varphi(z)$, $\deg\varphi = g \ge 1$. On the
$y$-class the edge operator is \emph{multiplication}:
$J(E, g_s) = [(ms + 1)\varphi + k z \varphi'] z^s$, and the top coefficient
$(mD + kg + 1)\varphi_g$ of the class ODE kills every polynomial candidate. Hence $E$ is
never the leading form of a Keller component, and $P = E + R$ (any strictly lower-weight
tail) is never a Keller component.
\end{theorem}

\begin{theorem}[{[MV/D]}; the vertex dichotomy (EXP-033)]\label{thm:t4}
If the monomial $x$ is a vertex of the Newton polygon $N(P)$ (convex hull of the support
with the origin) of a Keller component with linear part $x$, then $P = x + f(y)$.
\end{theorem}

\begin{lemma}[Annihilation (EXP-036)]\label{lem:annihilation}
For any polynomial $P$, $J(m, P^k) = -k\,L(P^{k-1} m)$ with $L = J(P, \cdot)$ (Leibniz). Any
covector annihilating the image of $L$ on a space containing the preimage therefore kills the
source; the certificate covectors used throughout are left-null on the completion matrix,
whose columns are exactly $L(m)$, so the tail steps of
Theorems~\ref{thm:t2}--\ref{thm:t4} are unconditional. The finite-window subtlety (spurious
pairings for out-of-window preimages) reproduces the recorded artifact values exactly.
\end{lemma}

The dichotomy's other side is real: $x + (x+y)^2$ and $x + (x+2y)^3$ are components in which
$x$ is \emph{swallowed} inside the polygon (equivalently: a pure power $x^d$, $d \ge 2$,
appears in the support). Every non-triangular component is of that kind. The residual open
core is therefore precise: components whose polygon swallows $x$ with a mixed staircase.

\section{The corner is diagonal; the staircase transport}\label{sec:staircase}

\begin{proposition}[{[MV]}; the diagonal corner (EXP-035)]
For $B = x^i y^j$, $i, j \ge 1$:
$J(B^p, x^a y^b) = p (ib - ja)\, x^{ip + a - 1} y^{jp + b - 1}$, with kernel exactly the
powers of $B$; no monomial gives $J(B^p, \cdot)$ a constant term. The mixed corner is clean;
the obstruction lives on the staircase of Newton vertices joining the swallowed linear
vertex to the corner.
\end{proposition}

\begin{theorem}[{[MV]}; block triangularity and the transport chain (EXP-037)]
\label{thm:transport}
Fix the $x$-edge grading $w = (v, 1-u)$ and let $s_0$ be the minimal $w$-class of $P$. In
the completion window (unknown $Q$ of bounded degree), every nonzero matrix entry sits at a
class offset $s - \sigma$ for a $P$-class $s$ ($\sigma = v + 1 - u$), and the offset
$s_0 - \sigma$ is minimal: ordering $Q$-classes by $w$-class, the system is block
triangular. Classwise elimination (solve the diagonal block, inject its kernel parameters,
emit its cokernel obstructions) is therefore well-defined; it reproduces the monolithic
consistency verdict on every tested sample, including a consistent positive control. On the
pure above-weight staircase family the diagonal blocks are the banded operators of
Theorem~\ref{thm:t1}, and on all $8$ grid samples the first inconsistent transport equation
sits at the \emph{constant's} class.
\end{theorem}

\begin{theorem}[{[MV]}; the fifth exclusion, window form (EXP-042)]\label{thm:t5}
Let $P = x + a\,x^u y^v + b\,x^d$ ($u \ge 2$, $v, d \ge 1$; the minimal shape swallowing
$x$). For every grid point $(u, v, d) \in \{2,3\} \times \{1,2\} \times \{2,3\}$ and window
$N = 7$, a cleared certificate covector with polynomial entries satisfies $c^T M = 0$
identically and pairs with the right-hand side to a \emph{monomial}:
\[
-13860\,a^7,\quad -32\,a^4,\quad -630\,a^4,\quad -2730\,a^7,\quad -1155\,a^3,\quad
-165\,a^4,\quad -4620\,a^7 b,\quad -81\,b^4,
\]
so the completion window is empty for \emph{all} parameter values with $a \ne 0$ (the
$b$-factors vanish only where Theorem~\ref{thm:t1} applies anyway). At $(u, v, d) =
(2, 1, 2)$ the pairing at window $N$ is $-c_N\,a^N$ with $c_N \in \{420, 1260, 13860,
180180, \dots\}$ for $N = 5..9$: a nonzero monomial in $a$ alone at every certified window.
The annihilation mechanism of Lemma~\ref{lem:annihilation} transfers verbatim (it is
$P$-agnostic), and in-window sources pair to zero exactly as the window criterion demands.
A three-parameter tail certificate ($-27720\,a^7 c_3^4$) extends the family.
\end{theorem}

\begin{lemma}[{[MV/D]}; the tower lemma (EXP-044/046)]\label{lem:tower}
Let $P = x + R$ ($R$ without constant or linear part), $M_N$ the completion-window
matrix at window $N \ge \deg P - 1$, and $T$ the total-degree top form of $P$. Suppose
every form in $\ker J(T, \cdot)$ at each degree is a scalar multiple of a power that
reduces, modulo $\C[P]$, to degree $\le N$ (this holds whenever $T$ is a primitive
monomial, and whenever $T = P - x$ or $T$ is the pure power $b x^d$, by expanding
$T^k = \lambda P^k + \text{lower}$). Then any left-null certificate of $M_N$ with
nonzero pairing extends to $M_{N+1}$ with the same (scaled) pairing: window
inconsistency propagates to every larger window. Proof ingredients, each
machine-verified: old columns vanish on new rows (degree bookkeeping); the new block is
$J(T, \cdot)$ on degree-$(N{+}1)$ forms with the predicted kernel; each resonance
$\kappa$ satisfies $\kappa - \lambda P^k \in$ window with $P^k \in \ker L$, so
$L(\kappa)$ is an in-window image which every certificate annihilates; the extension
was realized explicitly across an active resonance (window $8 \to 9$, ratio $2$,
pairing preserved). The measured window law ($-c_N a^N$ pairings with the odd-prime
$2N - 3$ ratios) is this tower seen through the clearing normalization.
\end{lemma}

\begin{theorem}[{[MV/D]}; the fifth exclusion, all degrees, primitive-top scope
(EXP-042/044/046)]\label{thm:t5all}
For $u \ge 2$, $v \ge 1$ with $x^u y^v$ primitive ($\gcd(u, v) = 1$) or $d \ge u + v$,
and any $a \ne 0$, $b$: $P = x + a\,x^u y^v + b\,x^d$ is not a Keller component at ANY
partner degree. (Base: the cleared certificates of Theorem~\ref{thm:t5}; induction:
Lemma~\ref{lem:tower}.) For proper-power tops (e.g. $u = v = 2$) the odd resonances do
not reduce mod $\C[P]$; the certificates kill them anyway in every tested instance
(e.g. $(xy)^5$ at window $9$), so the statement there is [D] pending the derivation of
that kill.
\end{theorem}

Two instrument-level results complete the toolkit. The \emph{matched-pair law} (EXP-038):
for cornered pairs $P = \alpha B^p + P_{\mathrm{sub}}$, $Q = \beta B^q + Q_{\mathrm{sub}}$,
the second class of $J(P, Q)$ is exactly
$\alpha p (ib - ja) q_{a,b} - \beta q (id - jc) p_{c,d}$ over outputs matched by
$(a, b) = (c, d) + (q - p)(i, j)$; imposing the pair shape does not change the obstruction
depth on any tested sample, but it halves the window unknowns: an efficiency tool, not a
second obstruction source. And the \emph{$x^m$-anchored edge operator} (EXP-043): for
$E = x^m \varphi(z)$, $z = x^k y^l$,
\[
J\left(x^m \varphi(z),\; x^{\alpha} y^{\beta}\right)
= x^{m + \alpha - 1}\, y^{\beta - 1}\,
\bigl[\beta m\, \varphi(z) + (\beta k - \alpha l)\, z\varphi'(z)\bigr],
\]
verified identically in symbolic edge coefficients; Theorem~\ref{thm:t3}'s operator is the
case $m = 1$.

\section{The frontier, recalibrated from primary sources}\label{sec:frontier}

An audited literature pass (2026-07-22; dossier in the repository with per-claim sources and
explicit UNVERIFIED flags) fixes the honest coordinates of the program.

\emph{Verified gcd coverage.} For a counterexample pair, $B := \gcd(\deg P, \deg Q)$
satisfies: $B \ge 9$ is forced by the classical interval results ($B = 1$: Magnus
\cite{Magnus1955}; $B \le 2$: Nakai--Baba; $B \le 8$ or prime: Appelgate--Onishi
\cite{AppelgateOnishi1985} and Nagata \cite{Nagata1989}); $B \ne p$ \cite{Abhyankar2008};
$B \ge 16$ (Heitmann \cite{Heitmann1990}); $B \ne 2p$ and $B = 16$ or $B > 20$
\cite{GGV2017}. Zoladek's claimed $B > 33$ has a documented gap \cite{GGV2017} and is not
used. Consequently our machine certificates at $\gcd = 2, 9, 12, 18$ (EXP-024..028), stated
in earlier revisions as frontier contact, are \emph{replications} of literature-covered
territory; we keep them as independent verifications and relabel them accordingly.

\emph{Verified degree floor.} Moh's classical analysis discards $\max\deg \le 100$
\cite{Moh1983}, with complete published detail only for the largest case ($64$); the
four exceptional cases correspond to six configurations in the modern corner terminology,
which the algorithmic audit of Guccione--Guccione--Horruitiner--Valqui enumerates and
discards \cite{GGHV2017, GGHV2022}. The current floor is: any counterexample has
$\max\deg \ge 125$, or $(\deg P, \deg Q) \in \{(72, 108), (108, 72)\}$ \cite{GGHV2022};
that single pair was left open explicitly for compute reasons.

\emph{The live targets.} (i) The $B = 16$ case comes with a forced direction set and
leading forms $l_{4,-1}(P) = \lambda_P R_0^m$, $R_0 = x (x y^4 - \lambda_0)^3$, and a
reduced normal form \cite{GGV2017}. We verify [MV] that $R_0$ is $(4,-1)$-homogeneous with
collinear $x$-anchored support, and, by the $x^m$-anchored operator above, that for
$m \ge 2$ \emph{no monomial reaches the Keller constant from the $R_0^m$ edge}: the
$B = 16$ problem is a staircase transport problem, the exact object of
Section~\ref{sec:staircase}. The pure $m = 1$ shape has $x$ as a Newton vertex and is
excluded outright by Theorem~\ref{thm:t4} (window replication empty); swallowed variants
land in Theorem~\ref{thm:t5}'s territory (sampled windows empty). (ii) The similarity filter
(partner supported in the scaled polygon $\tfrac{\deg Q}{\deg P} N^0(P)$; sound by the
similarity theorem since the scaled polygons are nested in the degree) cuts the
$(48, 64)$ full-window system from $2142$ to $111$ unknowns, and the validation runs are
DONE: three degree-48 samples of the $F_1$ flavor (the pure shape $x + R_0(1)^3$ and two
swallowed variants) have EMPTY filtered windows at $N = 64$ in $2.4$ to $3.6$ seconds
each, excluding all partners of degree $\le 64$ per sample (the classical case-64
degrees: a sampled-level replication of Moh/Heitmann demonstrating the pipeline's cost
at target scale). (iii) For the $(72, 108)$ system the histograms give $5992$ unknowns
in $535$ classes with largest diagonal block $22$ before filtering: with the filter and
the transport, the open territory (max degree $> 150$, and $(72, 108)$ itself) is
affordable at seconds-to-minutes per certificate.

\section{Positioning and novelty}\label{sec:positioning}

Per the adversarial novelty pass (same dossier): no prior statement of
Theorems~\ref{thm:t1}, \ref{thm:t2}, or of the companion's equivariant rigidity theorem was
found (for the latter, the $C^*$-equivariant analog \emph{fails} on pseudo-planes
\cite{DuboulozPalka2018}, and the finite-group analog is Miyanishi \cite{Miyanishi2023}, so
the $\C^2$ statement has content). Theorem~\ref{thm:t3} was not found as stated, but its
operator identity is within the standard bracket calculus \cite{Joseph1975, GGV2017}: the
statement appears new, the method is classical. Theorem~\ref{thm:t4}'s counterexample half
is essentially contained in \cite{GGV2017} (Prop.~4.1) with van den Essen's subrectangular
normalization \cite{vdEssen2000}; we claim only the sharp dichotomy form. All proofs are
elementary once seen; a folklore risk remains and is recorded, with the primary-source
audit trail (including its open verification items) persisted in the repository.

\section{Program}

Priorities, from the standing routes evaluation: (1) close the proper-power case of
Theorem~\ref{thm:t5all} (derive the odd-resonance kill measured in EXP-046), then extend
the tower induction to general above-weight tails and longer staircases; (2) transcribe
the $(72, 108)$ shape data and the beyond-150 $B = 16$ shapes and run the filtered
sweeps ($2$ to $4$ seconds per certificate at the validated scale); (3) the properness
reformulation ($JC(2)$ holds iff every planar Keller map has empty Jelonek set; exact
certificates shipped, EXP-014) as a cross-check instrument. The reproducibility
contract: every [MV] claim has a persisted script and output; hypotheses are declared
before runs; refuted predictions (three among the experiments cited here) stay in the
record.

\begin{thebibliography}{99}
\bibitem{Alpoge2026} L. Alp\"oge, announcement of the counterexample, X (@\_\_alpoge\_\_),
2026-07-19/20; map produced with Claude Fable 5 (Anthropic); question posed by Akhil.
\bibitem{Magnus1955} A. Magnus, On polynomial solutions of a differential equation,
\emph{Math. Scand.} \textbf{3} (1955), 255--260. (The gcd $=1$ theorem; Magnus' formula is
in \emph{Proc. Amer. Math. Soc.} \textbf{5} (1954), 256--266.)
\bibitem{AppelgateOnishi1985} H. Appelgate, H. Onishi, The Jacobian conjecture in two
variables, \emph{J. Pure Appl. Algebra} \textbf{37} (1985), 215--227.
\bibitem{Nagata1989} M. Nagata, Some remarks on the two-dimensional Jacobian conjecture,
\emph{Chinese J. Math.} \textbf{17} (1989), 1--7.
\bibitem{Moh1983} T.-T. Moh, On the Jacobian conjecture and the configurations of roots,
\emph{J. Reine Angew. Math.} \textbf{340} (1983), 140--212.
\bibitem{Heitmann1990} R. Heitmann, On the Jacobian conjecture, \emph{J. Pure Appl.
Algebra} \textbf{64} (1990), 35--72.
\bibitem{Abhyankar2008} S. S. Abhyankar, Some thoughts on the Jacobian conjecture, Parts
II--III, \emph{J. Algebra} \textbf{319} (2008), 1154--1248; \textbf{320} (2008),
2720--2826.
\bibitem{GGV2017} J. A. Guccione, J. J. Guccione, C. Valqui, On the shape of possible
counterexamples to the Jacobian conjecture, \emph{J. Algebra} \textbf{471} (2017), 13--74;
arXiv:1401.1784.
\bibitem{GGHV2017} J. A. Guccione, J. J. Guccione, R. Horruitiner, C. Valqui, Some
algorithms related to the Jacobian Conjecture, arXiv:1708.07936.
\bibitem{GGHV2022} J. A. Guccione, J. J. Guccione, R. Horruitiner, C. Valqui, Increasing
the degree of a possible counterexample to the Jacobian Conjecture from 100 to 108,
arXiv:2204.14178.
\bibitem{Joseph1975} A. Joseph, The Weyl algebra: semisimple and nilpotent elements,
\emph{Amer. J. Math.} \textbf{97} (1975), 597--615.
\bibitem{vdEssen2000} A. van den Essen, \emph{Polynomial Automorphisms and the Jacobian
Conjecture}, Progress in Mathematics 190, Birkh\"auser, 2000.
\bibitem{Miyanishi2023} M. Miyanishi, Equivariant Jacobian Conjecture in dimension two,
\emph{Transform. Groups} \textbf{28} (2023), 951--971.
\bibitem{DuboulozPalka2018} A. Dubouloz, K. Palka, The Jacobian Conjecture fails for
pseudo-planes, \emph{Adv. Math.} \textbf{339} (2018), 248--284.
\end{thebibliography}

\end{document}
